% Section 22.3: Ricci tensor and scalar curvature | Part: 4 | Depends on: sec:conn, sec:lc | Status: done | Original lines: 3208-3231
%!TEX root = ../../../main.tex
% REWRITTEN: sphere azimuth written \varphi instead of \phi (docs/NOTATION.md).
\subsection{Ricci tensor and scalar curvature}
\begin{formal}[{contractions of the curvature}]
The \emph{Ricci tensor} is the contraction of the contravariant index with the first covariant one, $\Ric=c(R)\in\TF{0}{2}(M)$:
\[
R_{jk}=\Ric_{jk}=R^i{}_{ijk}=\tr\big(v\mapsto R(v,\partial_j)\partial_k\big).
\]
When a metric $g$ is given (and $\nabla=\nabla^{LC}$), the \emph{scalar curvature} is the full contraction with the inverse metric
\[
s=g^{jk}\Ric_{jk}\in C^\infty(M).
\]
For $\nabla^{LC}$ the Ricci tensor is symmetric.
\end{formal}
\begin{intuition}[{Ricci = ``how volumes of small balls differ from flat ones''}]
$R$ has too many components to be pictured; Ricci averages it over directions. Roughly: a small ball of geodesics converging (positive Ricci) has less volume than in flat space; diverging (negative Ricci), more. In GR the Ricci tensor is sourced by matter: it is what appears in the Einstein equations; the remaining part of $R$ (the \emph{Weyl tensor}) describes gravitational waves and tidal fields in vacuum.
\end{intuition}

\begin{example}[{Ricci of the unit sphere}]
$R_{\varphi\varphi}=R^i{}_{i\varphi\varphi}=R^\theta{}_{\theta\varphi\varphi}+R^\varphi{}_{\varphi\varphi\varphi}=\sin^2\theta$ and $R_{\theta\theta}=R^\varphi{}_{\varphi\theta\theta}$. Compute the latter: $R^\varphi{}_{\varphi\theta\theta}=\partial_\varphi\Gamma^\varphi{}_{\theta\theta}-\partial_\theta\Gamma^\varphi{}_{\varphi\theta}+\Gamma^s{}_{\theta\theta}\Gamma^\varphi{}_{\varphi s}-\Gamma^s{}_{\varphi\theta}\Gamma^\varphi{}_{\theta s}=-\partial_\theta\cot\theta-\cot^2\theta=\csc^2\theta-\cot^2\theta=1$. Thus
\[
\Ric=\dd\theta\otimes\dd\theta+\sin^2\theta\,\dd\varphi\otimes\dd\varphi=g,\qquad s=g^{\theta\theta}\Ric_{\theta\theta}+g^{\varphi\varphi}\Ric_{\varphi\varphi}=1+1=2 .
\]
For the sphere of radius $a$, with its metric $g=a^2(\dd\theta^2+\sin^2\theta\,\dd\varphi^2)$: $\Ric=g/a^2$ (in 2D always $\Ric=\tfrac s2\,g$) and $s=2/a^2=2K$. (The components $R_{ij}$ do not depend on $a$, but $g$ and $g^{-1}$ do.)
\end{example}
